\subsection{带权强型不等式}
\label{sec:ch7-strong-weighted}

由 Marcinkiewicz 插值定理可从定理~\ref{thm:ch7-maximal-weighted-weak} 立即导出强型 $(p,p)$ 不等式. 设 $w\in A_q$, $1\le q<p$. 由 \eqref{eq:ch7-weight-set-comparison}, $w(E)=0$ 当且仅当 $|E|=0$, 所以 $L^\infty(w)=L^\infty$ 且范数相同. 因而
\[
 \|Mf\|_{L^\infty(w)}\le\|f\|_{L^\infty(w)}.
\]
在此不等式与弱型 $(q,q)$ 不等式之间插值, 得对所有 $p>q$,
\begin{equation}
\label{eq:ch7-maximal-weighted-strong}
 \int_{\mathbb R^n}|Mf|^pw\le C_p\int_{\mathbb R^n}|f|^pw.
\end{equation}
更深刻的结果是, 当 $p=q$ 时该不等式仍成立.

\begin{Theorem}
\label{thm:ch7-maximal-weighted-strong}
设 $1<p<\infty$. $M$ 在 $L^p(w)$ 上有界, 当且仅当 $w\in A_p$.
\end{Theorem}

强型不等式蕴含弱型不等式, 所以必要性由定理~\ref{thm:ch7-maximal-weighted-weak} 得到. 若能证明对给定 $w\in A_p$, 存在 $1<q<p$ 使 $w\in A_q$, 则充分性由上述插值论证得到. 这样的 $q$ 的存在性来自以下关键结果.

\begin{Theorem}[反向 Hölder 不等式]
\label{thm:ch7-reverse-holder}
设 $w\in A_p$, $1\le p<\infty$. 则存在只依赖于 $p$ 和 $w$ 的 $A_p$ 常数的常数 $C$ 与 $\varepsilon>0$, 使对任意立方体 $Q$,
\begin{equation}
\label{eq:ch7-reverse-holder}
 \left(\frac1{|Q|}\int_Qw^{1+\varepsilon}\right)^{1/(1+\varepsilon)}
 \le\frac C{|Q|}\int_Qw.
\end{equation}
\end{Theorem}

该名称来自如下事实: \eqref{eq:ch7-reverse-holder} 的反向不等式是 Hölder 不等式的直接推论. 为证明定理, 先证一个引理.

\hypertarget{chapter-07-page-150}{}
\begin{Lemma}
\label{lem:ch7-small-set-weight-decay}
设 $w\in A_p$, $1\le p<\infty$. 对每个 $0<\alpha<1$, 存在 $0<\beta<1$, 使对任意立方体 $Q$ 和 $S\subset Q$, 若 $|S|\le\alpha|Q|$, 则
\[
 w(S)\le\beta w(Q).
\]
\end{Lemma}

\begin{Proof}
在 \eqref{eq:ch7-weight-set-comparison} 中以 $Q\setminus S$ 代替 $S$, 得
\[
 w(Q)\left(1-\frac{|S|}{|Q|}\right)^p
 \le C(w(Q)-w(S)).
\]
若 $|S|\le\alpha|Q|$, 则
\[
 w(S)\le\left(1-\frac{(1-\alpha)^p}{C}\right)w(Q).
\]
故可取 $\beta=1-C^{-1}(1-\alpha)^p$.
\end{Proof}

\begin{Proof}[定理~\ref{thm:ch7-reverse-holder} 的证明]
固定立方体 $Q$, 在 $Q$ 内对 $w$ 作 Calderón--Zygmund 分解, 高度取递增序列
\[
 \lambda_0=\frac{w(Q)}{|Q|}<\lambda_1<\cdots<\lambda_k<\cdots,
\]
这些高度稍后确定. 对每个 $\lambda_k$, 得到一族两两不交的立方体 $\{Q_{k,j}\}$, 使
\[
 w(x)\le\lambda_k\quad\text{若 }x\notin\Omega_k=\bigcup_jQ_{k,j},
\]
且
\[
 \lambda_k<\frac1{|Q_{k,j}|}\int_{Q_{k,j}}w\le2^n\lambda_k.
\]
由构造显然有 $\Omega_{k+1}\subset\Omega_k$. 固定高度 $\lambda_k$ 分解中的 $Q_{k,j_0}$, 则 $Q_{k,j_0}\cap\Omega_{k+1}$ 是高度 $\lambda_{k+1}$ 分解中若干 $Q_{k+1,i}$ 的并. 因而
\[
\begin{aligned}
 |Q_{k,j_0}\cap\Omega_{k+1}|
 &=\sum_i|Q_{k+1,i}|\\
 &\le\frac1{\lambda_{k+1}}\int_{Q_{k,j_0}}w
 \le\frac{2^n\lambda_k}{\lambda_{k+1}}|Q_{k,j_0}|.
\end{aligned}
\]
固定 $\alpha<1$, 取 $\lambda_k=(2^n\alpha^{-1})^kw(Q)/|Q|$, 使 $2^n\lambda_k/\lambda_{k+1}=\alpha$. 于是
\[
 |Q_{k,j_0}\cap\Omega_{k+1}|\le\alpha|Q_{k,j_0}|.
\]
由引理~\ref{lem:ch7-small-set-weight-decay}, 存在 $\beta<1$ 使
\[
 w(Q_{k,j_0}\cap\Omega_{k+1})\le\beta w(Q_{k,j_0}).
\]

\hypertarget{chapter-07-page-151}{}
对高度 $\lambda_k$ 分解中的所有立方体求和, 得 $w(\Omega_{k+1})\le\beta w(\Omega_k)$. 迭代可得 $w(\Omega_k)\le\beta^kw(\Omega_0)$. 类似地, $|\Omega_k|\le\alpha^k|\Omega_0|$. 因而
\[
 |\Omega_\infty|=\left|\bigcap_k\Omega_k\right|=0.
\]
于是
\[
\begin{aligned}
 \frac1{|Q|}\int_Qw^{1+\varepsilon}
 &=\frac1{|Q|}\int_{Q\setminus\Omega_0}w^{1+\varepsilon}
 +\frac1{|Q|}\sum_{k=0}^\infty
 \int_{\Omega_k\setminus\Omega_{k+1}}w^{1+\varepsilon}\\
 &\le\frac{\lambda_0^\varepsilon w(Q)}{|Q|}
 +\frac1{|Q|}\sum_{k=0}^\infty\lambda_{k+1}^\varepsilon w(\Omega_k)\\
 &\le\frac{\lambda_0^\varepsilon w(Q)}{|Q|}
 +\frac1{|Q|}\sum_{k=0}^\infty
 (2^n\alpha^{-1})^{(k+1)\varepsilon}\lambda_0^\varepsilon\beta^kw(\Omega_0).
\end{aligned}
\]
取 $\varepsilon>0$ 使 $(2^n\alpha^{-1})^\varepsilon\beta<1$, 则级数收敛, 末项不超过 $C\lambda_0^\varepsilon w(Q)/|Q|$. 由于 $\lambda_0=w(Q)/|Q|$, 得到 \eqref{eq:ch7-reverse-holder}.
\end{Proof}

反向 Hölder 不等式给出完成定理~\ref{thm:ch7-maximal-weighted-strong} 所需的 $A_p$ 权性质.

\begin{Corollary}
\label{cor:ch7-ap-self-improvement}
有下列结论:
\begin{enumerate}[label=(\arabic*)]
\item $A_p=\bigcup_{q<p}A_q$, $1<p<\infty$;
\item 若 $w\in A_p$, $1\le p<\infty$, 则存在 $\varepsilon>0$ 使 $w^{1+\varepsilon}\in A_p$;
\item 若 $w\in A_p$, $1\le p<\infty$, 则存在 $\delta>0$, 使对任意立方体 $Q$ 和 $S\subset Q$,
\begin{equation}
\label{eq:ch7-ainfinity-condition}
 \frac{w(S)}{w(Q)}\le C\left(\frac{|S|}{|Q|}\right)^\delta.
\end{equation}
\end{enumerate}
\end{Corollary}

若权 $w$ 满足 \eqref{eq:ch7-ainfinity-condition}, 就说 $w\in A_\infty$. 这一记号来自第一条在 $p=\infty$ 时的形式.

\begin{Proof}
若 $w\in A_p$, 则由命题~\ref{prop:ch7-ap-properties}, $w^{1-p'}\in A_{p'}$. 因而它对某个 $\varepsilon>0$ 也满足反向 Hölder 不等式:
\begin{equation}
\label{eq:ch7-dual-reverse-holder}
 \left(\frac1{|Q|}\int_Qw^{(1-p')(1+\varepsilon)}\right)^{1/(1+\varepsilon)}
 \le\frac C{|Q|}\int_Qw^{1-p'}.
\end{equation}
取 $q$ 使 $q'-1=(p'-1)(1+\varepsilon)$, 则 $q<p$, 而 \eqref{eq:ch7-dual-reverse-holder} 与 $A_p$ 条件合起来推出 $w\in A_q$.

\hypertarget{chapter-07-page-152}{}
若 $p>1$, 取足够小的 $\varepsilon>0$, 使 $w$ 和 $w^{1-p'}$ 都满足指数 $1+\varepsilon$ 的反向 Hölder 不等式, 即得第二条. 若 $p=1$, 则对任意立方体 $Q$ 和几乎每个 $x\in Q$,
\[
 \frac1{|Q|}\int_Qw^{1+\varepsilon}
 \le C\left(\frac1{|Q|}\int_Qw\right)^{1+\varepsilon}
 \le Cw(x)^{1+\varepsilon}.
\]

为证第三条, 固定 $S\subset Q$, 并设 $w$ 满足指数 $1+\varepsilon$ 的反向 Hölder 不等式. 则
\[
 w(S)=\int_Q\chi_Sw
 \le\left(\int_Qw^{1+\varepsilon}\right)^{1/(1+\varepsilon)}
 |S|^{\varepsilon/(1+\varepsilon)}
 \le Cw(Q)\left(\frac{|S|}{|Q|}\right)^{\varepsilon/(1+\varepsilon)}.
\]
这给出 \eqref{eq:ch7-ainfinity-condition}, 其中 $\delta=\varepsilon/(1+\varepsilon)$.
\end{Proof}
